\subsection*{The complete appointed text}

The preceding formulary is part of this canonical full edition. Its Latin is
the image-collated text rather than normalized Bible Latin: the Introit retains
\latin{inimíci mei} and \latin{cógitant}; the Alleluia retains the comma after
\latin{Dómine}; the Gospel retains the comma after \latin{iumentum suum}; and
the Offertory and Communion preserve their received adaptations. The English
witnesses are printed without silently repairing those differences.

\elementguard
\subsection*{Introit: the whole Christ cries for immediate help \cue{Int.}}

Psalm 69 is the shortest complete prayer in the Psalter's Davidic collection.
Its title assigns it to David ``in remembrance,'' and its body repeats the
close of Psalm 39. Theodoret (PG~80, 1415--1418) places the prayer amid
Absalom's revolt: David remembers his own sin, sees danger surrounding him,
and hastens not to an earthly defense but to God. The literal voice is thus a
real servant in peril. The enemies seek his life and devise evil; he asks that
their attack collapse in shame and that all who seek God rejoice. The Mass
selects the opening crisis rather than the later turn to rejoicing. It brings
the congregation to worship as people whose first truth is need.

Eusebius (\work{Commentaria in Psalmos}, PG~23, 767--772) gives that need a
fully ecclesial range. The short prayer is to be kept ready on the lips; David
speaks, Christ speaks in the Passion, and Christ's members speak amid their
conflicts. Augustine, \work{Enarratio in Psalmum} 69.1--4, develops the same
mystery of the whole Christ. The Head who accepted human weakness and the body
still journeying through temptation do not pray as strangers to one another.
Their repeated \latin{Deus, in adiutórium meum inténde} confesses that grace is
not merely the beginning of salvation but the help by which every good act is
continued. Cassiodorus likewise receives the cry in Christ and the Church,
while St Robert Bellarmine notes both its Psalm 39 doublet and the Roman
Church's daily use.

The spiritual senses arise from that union. Allegorically, the suffering Head
teaches his members to seek the Father's help; morally, urgency becomes
watchful humility rather than panic; anagogically, the enemies' final confusion
and the rejoicing of those who seek God look toward deliverance completed.
Cassian, \work{Conference} 10.10, makes the verse a compendium for every
monastic state---assault, distraction, thanksgiving, compunction, and hope.
The Rule of St Benedict places it at the opening of the Hours and three times
on the incoming kitchen server's lips. Prayer orders both choir and labor.

The Introit therefore states the law of the entire formulary. The Collect's
runner begins by asking God to hasten. Paul's minister will confess sufficiency
from God. The wounded traveller will live because another approaches. Moses
will ask for a guilty people, and communicants will ask the mystery itself to
vivify and guard them. The cry is not left behind as the Mass proceeds; it is
answered in increasingly concrete forms of divine aid.

\elementguard
\subsection*{Collect: grace gives the service and enlarges the road \cue{Coll.}}

The Collect's grammar holds divine initiative and redeemed agency in one
sentence. In \latin{de cuius múnere venit, ut tibi a fidélibus tuis digne et
laudabíliter serviátur}, the first \latin{ut} clause names what comes from
God's gift. Yet \latin{a fidélibus tuis} names the faithful as real agents of
worthy and praiseworthy service. The imperative \latin{tríbue} then governs
\latin{nobis \ldots{} ut \ldots{} currámus}: God grants, and the faithful
themselves run. \latin{Ad promissiónes tuas} supplies the end; \latin{sine
offensióne} asks that obstacle, stumbling, and fall not interrupt the course.
The prayer is neither Pelagian self-motion nor passivity disguised as piety.
It asks for the freedom of children moved by grace toward their Father's word.

The appointed sequence makes the syntax audible. The Introit says,
\latin{Dómine, ad adiuvándum me festína}; the Collect answers,
\latin{tríbue \ldots{} ut \ldots{} currámus}. Paul then states the doctrinal
principle: \latin{sufficiéntia nostra ex Deo est}, and God
\latin{idóneos nos fecit minístros}. The Gradual enacts
\latin{laudabíliter}, because praise fills the mouth and the soul glories in
the Lord rather than in its course. In the Gospel, \latin{hoc fac, et vives}
and \latin{Vade, et tu fac simíliter} retain the full seriousness of action
after its source has been named. The lexical center is mercy: God is
\latin{miséricors}; the Samaritan is \latin{misericórdia motus}; the true
neighbor is the one who \latin{fecit misericórdiam}.

Augustine gives this order a rich biblical exposition. In
\work{Enarratio in Psalmum 118}, sermons 10.5--6 and 11.1--6, God grants what
unaided power cannot fulfill, enlarges the heart of the runner, brings the
killing letter to life by the Spirit, and orders love toward God and neighbor.
In \work{De gratia et libero arbitrio} 16.32--17.33 he insists that the divine
command summons a real human will precisely because grace heals and assists
it. Hilary's treatment of Psalm 118:32, \work{Tractatus in Psalmum 118},
Daleth 41--42 (PL~9, 303--304), likewise joins running to the widened way of
the commandments. Benedict's Prologue turns the same movement into a school
of service whose narrow beginning opens into running with an expanded heart.

Aquinas gathers the doctrine in \work{ST} I--II, q.~109, aa.~2, 4--5, 9--10:
grace heals the will, makes the good act possible in its supernatural order,
enables fulfillment of the law, and sustains perseverance. Trent, Session VI,
chs.~5--6 and 16,
teaches that preparation is impossible without prevenient grace, that the
moved person freely assents and cooperates, and that growth in justice is at
once God's gift and the justified person's living action. The Collect prays
this doctrine before it defines it.

The witnesses attest the prayer in differing forms. A shorter cognate form in
the Verona collection (Feltoe, p.~74) omits
\latin{quaesumus nobis}. The Old Gelasian places the second Collect in Book
III, no.~VIII (Wilson, pp.~228--229), with an apparatus variant
\latin{a promissionibus}. The Hadrianum gathers the three orations under
\emph{Dominica XIII post Pentecosten}, no.~XXXI (Wilson, p.~173, Cambrai 164,
fol.~167v). These identified witnesses display differing forms and the
assembled oration set without proving a direct line of descent.
Schuster (III, pp.~128--129) reads the prayer directly with 2~Cor.~3:5:
God-given sufficiency is the power by which the worshipper runs.

\elementguard
\subsection*{Epistle: the Spirit makes ministers of the new covenant \cue{Ep.}}

Second Corinthians 3 is not an abstract tract on texts. Paul has called the
Corinthian Church his letter of commendation, written by Christ through
apostolic ministry, not with ink but with the Spirit, not on stone but on
hearts (vv.~1--3). He then locates apostolic confidence
\latin{per Christum ad Deum}. Neither natural talent nor self-certification
makes the minister sufficient; God makes him fit for a new-covenant ministry.
The contrast with stone and death invokes Moses' shining face in Exod.~34 and
the promises of a heart-written covenant in Jer.~31:31--34 and
Ezek.~36:25--27.

The appointed ending at v.~9 is important. Paul has called the old ministry
glorious and the new ministry more glorious; vv.~10--18 continue the argument
through surpassing glory, the veil, conversion to the Lord, freedom in the
Spirit, and transformation into Christ's image. Within that whole,
\latin{líttera occídit} cannot make Scripture's literal sense, the moral law,
Moses, or Israel evil. The Law is God's gift and truly reveals his will. As
Rom.~7--8 explains, command encountered externally by the sinner exposes and
condemns but does not of itself forgive or pour charity into the heart. The
Spirit's ministry gives righteousness and the power of a filial obedience.

St John Chrysostom, \work{Homilies on 2 Corinthians} 6.2 and 7.1, 5, begins
where Paul begins: the Corinthians are the apostolic letter because the Spirit
has truly formed them. Confidence is not vanity but trust in God, and Moses'
ministry possesses a glory that the new gift surpasses rather than falsifies.
Theodoret (PG~82, 391--394) reads sufficiency as a divine making-fit and the
Spirit as the power by which obedience becomes possible. Ambrosiaster
(PL~17, 176--177) finds the contrast in effects: the letter condemns sin,
whereas the Spirit gives forgiveness and righteousness. These witnesses honor
the first dispensation while proclaiming fulfillment.

Augustine builds his mature theology of grace around this passage. In
\work{De spiritu et littera} 4.6, 14.23--25, 19.34, and 26.42, the killing
letter is the good command left external to a heart still ruled by
concupiscence; the vivifying Spirit writes love within, so that righteousness
is not only known but desired and done. His earlier \work{De doctrina
christiana} 3.5.9 applies the phrase to carnal captivity to signs, a genuine
hermeneutical extension but not the primary covenantal sense. Aquinas,
\work{Super II ad Corinthios} 3, lects.~1--2, keeps divine causality and human
ministry together, while \work{ST} I--II, q.~106, aa.~1--2 identifies the New
Law principally with the Holy Spirit's grace. Even the written Gospel would
remain external letter if that grace were absent.

The allegorical fulfillment is Christological: in the incarnate and risen
Lord the promised covenant is enacted and the Spirit is poured out. Morally,
the Epistle forbids both boasting in talent and hiding from service; God-given
fitness exists for ministry. Anagogically, the movement from glory to
surpassing glory opens beyond the pericope toward transformation into the
Lord's image. Ritually, the reading prepares the Gospel's demonstration that
the Spirit does not produce a vague interiority. He makes the written love of
God and neighbor become merciful deed, and he makes a restored traveller into
one capable of bearing another.

Catholic doctrine therefore reads Paul within one economy of revelation.
\work{Dei Verbum} 14--16 honors the Old Testament's lasting value and
Christological fulfillment; \work{Nostra aetate} 4 and the Pontifical Biblical
Commission's \work{Jewish People and Their Sacred Scriptures} 41 forbid using
2~Cor.~3 as a judgment on Jewish persons or as a repudiation of Israel's
Scriptures. The contrast is between ministries and modes under sin and grace,
not between peoples ranked by worth.

\elementguard
\subsection*{Gradual: David's praise becomes the Church's continual boast \cue{Grad.}}

Psalm 33 bears a received Davidic title: David changed his countenance before
the Philistine ruler and escaped. First Samuel names Achis, while the title
names Abimelech. Basil, Augustine, and Theodoret preserve different solutions,
including a dynastic title and a spiritual naming. The historical sense
remains David's praise after deliverance during Saul's persecution. The
alphabetic form gives the thanksgiving a wisdom-like completeness: the rescued
man teaches the meek how fear, praise, speech, conduct, and trust belong
together.

St Basil, \work{Homilia in Psalmum 33} 1 (PG~29, 352--357), takes
\latin{in omni témpore} with its full moral force. Praise is not confined to
prosperity; Job's blessing under loss displays a mouth disciplined for every
condition. Eusebius (PG~23, 288--297) moves from Davidic thanksgiving toward
the Church assembled in the Lord. Theodoret (PG~80, 1101--1109) preserves the
flight and rescue while teaching praise amid pursuit and calm. Bellarmine
applies ``always'' to fitting praise throughout life and to the Church's
unceasing public worship. The meek rejoice because the psalmist's boast is not
private greatness but the Lord.

Augustine preached Psalm 33 in two sermons. He reads its Abimelech title
alongside an Old-Latin/Septuagint-shaped narrative lemma from 1~Samuel 21:13,
``he was carried in his own hands,'' and applies the phrase to Christ at the
Supper, carrying in his hands the sacrament of his own Body. That clause
belongs to Augustine's Samuel text, not to the psalm title; it is absent from
the Clementine Samuel and from the appointed Gradual. The juxtaposition
nevertheless generates a profound Eucharistic reception:
David's deliverance-song becomes the voice of the true David who gives himself
as food. The literal rescue and sacramental mystery are ordered rather than
competitive. The delivered king praises; Christ gives the definitive
deliverance; the Church about to offer and communicate learns to glory in him.

Within the Mass the Gradual is the answer to the Epistle. Paul refuses
\latin{quasi ex nobis}; the psalmist answers, \latin{in Dómino laudábitur
ánima mea}. God has made ministers sufficient, and their praise returns the
gift to its source. The verse also prepares the lawyer's question. Eternal life
will not be possessed by a self-justifying boast, but by love whose glory is
God and whose joy becomes shared with the meek.

\elementguard
\subsection*{Alleluia: Christ and his members pray through the night \cue{All.}}

Psalm 87 begins by naming God as salvation and then descends through affliction,
abandonment, the pit, wrath, and darkness. Its received heading assigns the
instruction to Heman the Ezrahite, or in another Catholic line to David
speaking in his name. Haydock relates Heman to the early monarchy and presents
the whole psalm as agreeing to Christ in his Passion, death, and burial. The
first verse already contains the mystery: prayer continues by day and by
night, not because sorrow is denied but because the afflicted one still says
\latin{Deus salútis meæ}.

Augustine, \work{Enarratio in Psalmum} 87.1--2, hears Christ speak in the
Passion and his body pray in him. Day and night also signify prosperity and
adversity; the Church offers one prayer through both. Cassiodorus continues
this Christological-ecclesial reading, and Bellarmine places its historical
fullness in Gethsemane and Good Friday. The darkness is therefore taken into
the Son's obedient prayer. It is not permitted to define the final relation
between suffering and God apart from the crucified Christ.

Eusebius (PG~23, 1052--1056) reads Psalm 86's proclamation of birth in Sion
followed by Psalm 87's descent toward death as an incarnational sequence. Day
and night become the incarnate Lord's continual prayer and oblation. Theodoret
(PG~80, 1567--1573) hears Israel in calamity and, through Israel, the sorrows
of fallen humanity. These senses enlarge the Church's voice: Christ prays,
Israel laments, humanity cries, and the baptized pray as Christ's members.

The choice of this verse for an Alleluia is a liturgical act of hope. The
Gradual's \latin{in omni témpore} is not tested only by alternating good and
bad moods; it enters the Lord's own night. Immediately afterward Jesus calls
the disciples' eyes blessed. The juxtaposition teaches that revelation is
gift, while faith's prayer continues when consolation is not seen. The
anagogical sense lies already in the address: the God invoked as salvation
will not let the night possess the final word.

\elementguard
\subsection*{Gospel: Christ the Neighbor creates neighbors by mercy \cue{Gosp.}}

Luke places the episode within Jesus' journey toward Jerusalem. Christ first
turns privately to the disciples and blesses the eyes that see what prophets
and kings desired. A lawyer then tests him by asking what deed leads to eternal
life. Jesus sends the expert to the Law: \latin{Quid scriptum est? quómodo
legis?} The lawyer joins Deut.~6:5 and Lev.~19:18, loving God wholly and the
neighbor as oneself. Christ confirms the answer: \latin{Hoc fac, et vives}.
Revelation, reading, and life belong together.

The second question discloses the resistance: \latin{volens iustificáre
seípsum}, the lawyer asks for the neighbor's boundary. Jesus answers by
narrating embodied compassion. A man descends from Jerusalem to Jericho and is
left half-dead. Priest, Levite, and Samaritan each see; what distinguishes the
Samaritan is the movement after sight. He is moved with compassion, approaches,
binds the wounds, pours oil and wine, places the man on his own beast, brings
him to an inn, tends him through the night, pays two denarii, entrusts further
care, and promises repayment upon return. Jesus reverses the grammar: not
``which person counts as my neighbor?'' but ``which became neighbor?'' The
answer is the one who did mercy; the command is \latin{Vade, et tu fac
simíliter}.

The Jerusalem--Jericho descent and its desert terrain make the danger
concrete. Luke gives a parable and does not identify the victim; Cornelius a
Lapide judges it drawn from a frequent real occurrence and calls it a true
story in that sense. The evangelist does not state why priest and Levite pass
by. The Samaritan is the distrusted outsider who fulfills the Law's command,
so the narrative rebukes every attempt to confine mercy by status. It cannot
be turned into contempt for Jewish priesthood or people: Jesus, the lawyer,
the double commandment, and the story's prophetic world all stand within
Israel's revelation.

The ancient spiritual reading contemplates the complete economy of salvation
within the story. Origen's \work{Homily on Luke} 34.3 reports an elder's
allegory, which Ambrose VII.73--84, Augustine
\work{Quaestiones Evangeliorum} II.19, and Bede further develop. Adam and
humanity descend from the heavenly Jerusalem into Jericho, the mutable mortal
world. The devil and his ministers strip the robe of grace and immortality and
wound through sin. Humanity remains half-alive: rational nature and knowledge
of God are not annihilated, yet the power to attain the saving end has been
grievously wounded. Priest and Levite figure an Old Covenant ministry that is
holy but cannot of itself heal fallen nature. The Samaritan---the Guardian---is
Christ, who becomes near by assuming our condition.

Each act then discloses the Physician's mercy. Christ binds the disorder of
sin; in Ambrose, oil signifies remission and wine the announcement of judgment
that brings compunction; in Augustine, oil is the consolation of good hope
through pardon and reconciliation, wine the exhortation to fervent action.
His beast is the flesh of the Incarnation that bears our weakness. The inn is
the Church, where the wounded are restored through the ministry Christ
entrusts. The two denarii signify the two commandments of charity, or the
promises of present and future life, in Augustine, and the two Testaments in
Ambrose. The morrow opens toward the
Resurrection; the innkeeper figures apostolic and pastoral care; the promised
return is the Parousia and judgment, when faithful expenditure for the wounded
is rewarded. Origen 34.3 explicitly attributes the return-reading to the elder
whose interpretation he transmits.

This allegory is not an ornamental afterthought. It answers the deepest
question raised by the moral command: who first becomes neighbor to fallen
humanity, and whence comes power to do likewise? Christ approaches before the
wounded man can move. He bears, pays, entrusts, and returns. The moral sense
therefore flowers from the allegorical: receive the Lord's mercy, then bind
another's wounds, bear burdens, spend material goods, provide durable care,
and exercise charity under prudence and justice. Bede states the union
positively: everyone who shows mercy is neighbor, and Christ is Neighbor above
all. The spiritual sense intensifies rather than attenuates the institution of
mutual brotherhood.

The tradition develops that union in complementary registers. Clement of
Alexandria, \work{Quis dives salvetur?} 28--29, presents Christ as the supreme
neighbor whose love generates shared love. Irenaeus, \work{Against Heresies}
III.17.3, shows Christ compassionating the man fallen among thieves, binding
his wounds, entrusting him to the Holy Spirit, and giving the two royal
denarii. Cyril of
Alexandria, Sermons 67--68, emphasizes the fulfillment of the Law and Christ's
pedagogical reversal of the lawyer. Theophylact (PG~123, 843--852) joins
neighborly action to the Christological Samaritan. Albert (Borgnet XXIII,
55--70) enumerates the Samaritan's seven acts before opening the spiritual
senses. Bonaventure (Quaracchi VII, 266--272) develops road, robbery, clerical
duty, remedies, universal neighbor-love, and Jesus' pedagogy within a single
commentary.

Anagogically, the lawyer's first desire is not discarded: eternal life is the
end. The pilgrim Church is the inn where humanity is healed on the way to the
heavenly patria; the Resurrection is the morrow; the Lord will return and
bring hidden expenditure into judgment. At this Mass the allegory becomes
ritual action. The Collect asks to run toward the promises; the Epistle names
the Spirit who vivifies; the Offertory will show the mediator before God; the
Secret asks pardon; Communion and Postcommunion give the life by which the
healed go and act. St John Paul II's \work{Salvifici doloris} 28--30 and
\work{Dives in misericordia} 3--4, Benedict XVI's \work{Deus caritas est}
15 and 20--25, and Francis's \work{Fratelli tutti} 56--86 continue the
patristic union of received divine love and concrete service to the wounded.

\elementguard
\subsection*{Offertory: Moses' intercession enters Christ's self-offering \cue{Off.}}

Exodus 32 must be heard across the covenant narrative. Israel has received
deliverance, the commandments, blood of covenant, and the promise of God's
presence. While Moses remains on the mountain, the people fashion the calf.
God reveals the apostasy to Moses and announces judgment; the very revelation
raises the intercessor. Moses appeals not to Israel's innocence but to the
divine work already begun, God's name among the nations, and the oath to
Abraham, Isaac, and Jacob. Mercy preserves the people so that covenantal
presence may be restored in Exodus 33--34, although sin still bears judgment
and discipline.

The Offertory's received Latin is an ancient witness in its own right. It
condenses Exod.~32:11--14, incorporates matter from v.~12, names Jacob where
the Clementine Vulgate has Israel, substitutes the land flowing with milk and
honey for the stars-and-land wording, and ends with
\latin{placátus factus est Dóminus de malignitáte, quam dixit fácere pópulo
suo}. Augustine's Old-Latin lemma in \work{Quaestiones in Heptateuchum}
II.143 independently transmits the same construction. The Church therefore
sings the text she received rather than silently normalizing it to the later
Vulgate.

Patristic exegesis enters the divine-human drama without reducing it. In
questions II.143, 147, and 149--151 Augustine explains that the threatened
``evil'' is penal judgment, never moral evil in God; the eternal will includes
the temporal signs and creaturely prayers through which providence works. His
\work{City of God} XV.25 names divine anger as just action, not a perturbation
of the immutable divine life. Theodoret, \work{Questions on Exodus} 67--68,
likewise reads the anthropomorphic language according to divine constancy.
Question 68 preserves the moral seriousness of the sequence: mercy delays and
preserves, while correction and accountability remain.

The Fathers then see Moses as pastor and type. First Clement 53 remembers the
servant willing to be blotted from God's book for the people. Chrysostom,
\work{On the Priesthood} IV.1, finds there the daring love required of a
pastor; his \work{Homily 16 on Romans} sees Moses risk everything for God's
glory and the salvation of those entrusted to him. Gregory the Great,
\work{Moralia} IX.16.23, reads ``let me alone'' as God himself giving the
intercessor confidence to plead; \work{Moralia} XX.5.14 treats Moses'
willingness to perish with the people as self-offering. Tertullian,
\work{Adversus Marcionem} II.26, explicitly names Moses' deprecation a type of
Christ.

The type reaches its truth at the altar. Moses stands before God for the
people; Christ is the unique Mediator who gives himself for sinners and whose
intercession does not cease. Aquinas, \work{ST} I, q.~19, a.~7, ad~1 and
II--II, q.~83, a.~2, explains how immutable providence wills prayer itself as
a real secondary cause. The \work{Catechism} 2574--2577 presents Moses as
formed in faithful, persevering intercession, while Benedict XVI's audience of
1 June 2011 places his plea within the mystery fulfilled in Christ.

The complete formulary makes the type fruitful in four senses. Literally,
Israel is preserved through the mediator God raises. Allegorically, Moses
figures Christ who stands for sinners and offers himself. Morally, the minister
made sufficient by God cannot advance by abandoning the fallen; he enters the
need of the people in prayer. Anagogically, the patriarchal inheritance and
restored covenant anticipate the promised homeland and the final healing for
which the Church runs. At the Offertory the congregation places gifts upon the
altar while this intercession is sung: judgment is confessed, promise is
remembered, and pardon is sought within Christ's one propitiatory oblation.

\elementguard
\subsection*{Secret: the true propitiatory offering glorifies God's name \cue{Sec.}}

The Secret asks God to look favorably upon the offerings:
\latin{Hóstias, quǽsumus, Dómine, propítius inténde}. The ensuing
\latin{ut} clause gives their Godward end. The grammatical subject of
\latin{dent} is still \latin{hóstiae}; \latin{nobis indulgéntiam largiéndo}
describes how they give honor to God's name---by granting pardon to us.
\latin{Largiendo}, an ablative gerund, does not grammatically make God its
immediate subject, although all efficacy and pardon come from him. Cummiskey's 1861
English changes both order and agency, so it remains a historical translation,
not a syntactic control for the Latin.

The prayer is attested in differing forms. The Old Gelasian Book
III, no.~VIII joins the three orations and records the
\latin{largiaris}/\latin{largiendo} variation. The Hadrianum's
\emph{Dominica XIII post Pentecosten}, no.~XXXI (Wilson, p.~173), gives the
later set with \latin{largiendo}; the 1962 Missal retains that reading.
Schuster (III, pp.~128--131, especially p.~131) reads indulgence as the
central favor to which the formulary's benefits are ordered. In its received
form, the prayer offers a compact sacrificial theology: the Church offers, God
regards propitiously, pardon is truly granted, and the divine name receives
honor.

The ancient witnesses explain how creaturely offering belongs within Christ's
gift. Justin, \work{Dialogue with Trypho} 116--117, sees Malachi's pure
offering fulfilled in the Eucharistic prayers and thanksgiving of Christians.
Irenaeus, \work{Against Heresies} IV.17.5--18.6, insists that the Church offers
created bread and wine in freedom and gratitude, confessing the Creator while
receiving the firstfruits sanctified. Cyril of Jerusalem,
\work{Mystagogical Catechesis} V.8--10, presents the Eucharistic sacrifice as
the memorial of Christ's offering within which the Church intercedes for all.
Augustine, \work{City of God} X.5--6, 20, defines true sacrifice as the work
that unites us to God and finds its perfect form in Christ, priest and offering,
with the Church his body offered in him.

Trent, Session XXII, ch.~2, gives the prayer's doctrinal register: the Mass is
a true and proper propitiatory sacrifice. The victim and priest are the same
Christ as on Calvary; the manner of offering is sacramental and unbloody; by
this oblation God truly grants grace and the gift of repentance and pardons
crimes and sins. Therefore the Secret's positive claim must be allowed its
full weight. These offerings, taken into Christ's action, really are the
Church's sacrificial means of asking pardon and honoring the Father. They do
not rival the Cross or add another redemption; they make sacramentally present
the one sacrifice from which every pardon flows.

The cross-proper sequence is now clear. Moses has pleaded a prior promise for
a guilty people. At the altar, his typical intercession enters the prayer of
Christ and his body. The self-justifying lawyer has been answered by mercy; the
Secret teaches the worshipper to seek justification as pardon before offering
obedience as glory. And because pardon gives honor to God's name, mercy does
not terminate in the recipient. It returns as praise, sacrificial
thanksgiving, and the neighbor-love commanded in the Gospel.

\elementguard
\subsection*{Communion: providence flowers into sacramental abundance \cue{Comm.}}

Psalm 103 is a vast hymn to the Creator who orders the heavens, waters,
mountains, animals, seasons, labor, and sea. Verses 13--15 belong to its
providential center: God waters the mountains from his upper chambers; earth is
filled from his works; grass serves animals and plants serve human cultivation;
bread strengthens the heart, wine gladdens it, and oil brightens the face. The
Latin psalter receives the psalm under David's name. The literal sense is not
generic ``nature'' but creation continually receiving being, order, fertility,
and usefulness from the Lord.

The Communion antiphon makes a precise liturgical selection. It begins with
the second half of v.~13, reorders the clause, inserts \latin{Dómine}, omits
the first half of v.~14, and continues with the bread, wine, and oil of v.~15.
The resulting first-person address turns cosmological praise into the
communicants' thanksgiving: the earth is filled from \emph{your} works, Lord.
Nothing created is despised at the moment of sacramental reception. The fruits
of earth and human labor remain within the Creator's economy and are brought
to their highest liturgical service.

The Greek tradition first unfolds the providential abundance. Theodoret
(PG~80, 1693--1707) explains the bodily gifts directly: bread sustains, wine
brings fitting gladness, and oil nourishes and brightens. His reading gives
Christian asceticism its grateful basis; temperance is ordered use, not
contempt for the Creator's goods. Eusebius (PG~23, 1273--1280) begins from the
same providence and opens spiritual senses: wine calls forth the true Vine,
bread the mystical Logos, oil illumination and healing. Created signs become
capable of speaking Christ because all things were made through the Word.

Augustine, \work{Enarratio in Psalmum} 103, sermo 3.12--14, draws the passage
into the whole Christ. Apostolic earth bears the bread of Christ preached to
the nations; the cup gives a sober inebriation that lifts the heart without
destroying reason; oil signifies the anointing that runs from the Head through
his members. Cassiodorus, \work{Expositio Psalmorum} 103.20 (PL~70,
733--734), gives the ecclesial mystery explicit sacramental form: wine is
consecrated as Christ's Blood, and oil belongs to royal chrism. Bellarmine
completes the moral reception with gratitude, recognition of fatherly
providence, and temperance.

At Communion these senses stand in their Catholic order. Rain, cultivation,
bread, wine, and oil are truly good; their providential service prepares the
sacramental economy. The Church now receives not simply a symbol of absent
care but Christ himself under the consecrated species. The Psalm's abundance
becomes thanksgiving for a superabundant gift: the Creator has taken what
creation and labor supplied and, through Christ's words and sacrifice, gives
the faithful his Body and Blood. The Gradual's Eucharistic reception of the
true David and the Secret's theology of offering meet here in sacramental
possession.

\elementguard
\subsection*{Postcommunion: holy participation vivifies, expiates, and guards \cue{Postcomm.}}

The Postcommunion speaks after reception and still asks. Its subject,
\latin{huius participátio sancta mystérii}, is holy participation in this
mystery, not a merely subjective recollection. The first petition,
\latin{vivíficet}, gathers Paul's \latin{spíritus vivíficat}; the second asks
that the same participation grant both \latin{expiatiónem} and
\latin{munímen}. \latin{Páriter} coordinates the two gifts: cleansing from
sin and defense for the continuing contest belong to one sacramental grace.
The Church asks for objective divine action to bear its proper fruit in those
who have communicated.

The text preserves a long prayer-family. The Verona collection has
\latin{participatio tui sancta mysterii} (Feltoe, p.~60); the Old Gelasian and
the Hadrianum no.~XXXI transmit the prayer within the three-oration set; the
1962 form reads \latin{huius}. The change from \latin{tui} to \latin{huius}
shifts the demonstrative reference but not the theological structure.
Participation remains a gift whose life, expiation, and protection are sought
from God.

St Ignatius of Antioch, \work{Letter to the Ephesians} 20, calls Eucharistic
communion the medicine of immortality and antidote against death, because it
is union with Jesus Christ. Aquinas, \work{ST} III, q.~79, aa.~1--2, 4, 6, 8,
distinguishes the effects contained in this compact prayer. The sacrament gives
spiritual life and growth as food; it remits venial sin and strengthens
charity; it preserves from future sin by union with Christ and the spiritual
delight that resists disordered desire; and it pledges glory. These effects
flow from Christ's Passion made sacramentally present, while fruitful reception
requires living faith and a fitting disposition.
The \work{Catechism} 1391--1395 receives the same doctrine in pastoral form.

The final prayer consequently seals every movement of the Mass. The Introit's
helpless cry receives living aid. The Collect's runner is nourished and
defended on the road. The Epistle's Spirit-given life takes sacramental form.
The Gradual and Alleluia's continual praise is sustained through trial. The
Gospel's wounded humanity is not only commanded but healed in Christ, and its
new life is strengthened for merciful action. Moses' plea and the Secret's
propitiatory petition reach pardon through the one sacrifice. Communion's
created abundance reaches Christ's gift of himself. The communicant is thus
sent out neither self-sufficient nor empty-handed: vivified, expiated, and
guarded, he can go and do likewise while running toward the promises.
